\begin{definition}[The accessible two-register charge-pair operator]
    \label{def:peps_regular_charge_pair}
    \lean{TNLean.PEPS.regularChargePairCoefficient,
        TNLean.PEPS.regularChargePairOperator,
        TNLean.PEPS.normalizedRegularChargePairCoefficient}
    \leanok
    Let $G$ be a finite group, $\chi:G\to\mathbb C$, and $p\in G$.
    On the two group registers, define the diagonal operator
    \[
        \Pi(\chi,p)|g,h\rangle=\chi(ph^{-1}g)|g,h\rangle.
    \]
    Its coefficient vector is $v_{\chi,p}(g,h)=\chi(ph^{-1}g)$;
    put $\widehat v_{\chi,p}=|G|^{-1}v_{\chi,p}$.
    This is the accessible operator of
    \cite[charge-pair creation]{Schuch2010PEPS}, lines 2505--2535.
    The identities below concern the two registers; an original-spin
    creation operation requires a separate physical contraction comparison.
\end{definition}

\begin{theorem}[Simultaneous-translation invariance of the actual charge pair]
    \label{thm:peps_regular_charge_pair_invariance}
    \lean{TNLean.PEPS.regularChargePairCoefficient_left_invariant,
        TNLean.PEPS.regularChargePairOperator_commute,
        TNLean.PEPS.regularChargePairOperator_conjugation}
    \leanok
    \uses{def:peps_regular_charge_pair}
    For every $x,g,h\in G$ and every coefficient function $\chi$,
    \[
        v_{\chi,p}(xg,xh)=v_{\chi,p}(g,h).
    \]
    If $U_x$ is the left-regular permutation matrix, then
    \[
        [U_x\otimes U_x,\Pi(\chi,p)]=0,
        \qquad
        (U_x\otimes U_x)\Pi(\chi,p)(U_x\otimes U_x)^*
        =\Pi(\chi,p).
    \]
    No invariance premise on $\chi$ is needed.
\end{theorem}
\begin{proof}\leanok
    The relative coordinate satisfies $(xh)^{-1}(xg)=h^{-1}g$.
    Thus the diagonal entries are unchanged by the simultaneous permutation.
    The permutation matrix is unitary, giving the conjugation identity.
\end{proof}

\begin{theorem}[Normalization and the nontrivial-charge orthogonality]
    \label{thm:peps_regular_charge_pair_normalization}
    \lean{TNLean.PEPS.regularChargePairCoefficient_normSq,
        TNLean.PEPS.normalizedRegularChargePairCoefficient_normSq,
        TNLean.PEPS.normalizedRegularChargePairCoefficient_orthogonal_uniform}
    \leanok
    \uses{def:peps_regular_charge_pair,
        thm:peps_translated_character_inner_product}
    If $\chi$ is the character of a unitary irreducible representation, then
    \[
        \|v_{\chi,p}\|^2=|G|^2,
        \qquad \|\widehat v_{\chi,p}\|^2=1.
    \]
    If this irreducible character is nontrivial, the normalized vector is
    orthogonal to the normalized uniform vector
    $\Omega(g,h)=|G|^{-1}$. This latter orthogonality needs only
    irreducibility and nontriviality, without unitarity of the representation.
\end{theorem}
\begin{proof}\leanok
    For each fixed $h$, translated character orthogonality gives squared
    norm $|G|$ upon summing over $g$. Summing over $h$ gives $|G|^2$.
    The nontrivial character has zero pairing with the trivial character
    for every translation, so its pairing with the uniform pair vector vanishes.
\end{proof}
